% Pursuing Stacks (cotes 134-2 to 134-8): archive cote 134-3, site page 50 (PDF page 51), typescript page 50.
% typescript. Draft compared: ps2.tex:2488-2564.
% AI-transcribed by Claude Opus 5.5 (claude-opus-5-5, adaptive thinking, effort high), Message Batches API, batch ps1 (msgbatch_01Rv5hCaNTtUt4JHx1yqoHPL), 2026-11-12; not yet checked by eye.
\scanpage{134-3/50}
\mspage{50}

Now let's say that $e_{A^\wedge}$ is a \xout{homotx} \uline{deformation retract of} F (F any ele-%
ment in $A^\wedge$) if the identity map of F is homotopic (with respect to I) to
a``constant''map of F into itself.  It is purely formal,  using the diagonal
map $I\to I\mathrm{x}I$ in $A^\wedge$,  that if (more generally) $F_o$ is a deformation re\xout{x}tract
of F,  and $G_o$ a deformation retract of G ,  then $F_o\mathrm{x}G_o$ is a deformation
retract of FxG. Thus, a \uline{sufficient} condition for T 2) to hold is the fol-%
lowing ``homotopy\inserted{-test} axiom'' :

\uline{Condition} \xout{\illegible{1}} T H :  For any  a in A,  \xout{\unsure{the}}\inserted{e=} $e_{A^\wedge}$ is a deformation retract of
a with respect to $(I,e_o,e_1)$,  namely there axists a
(1)  section  $c_a : e\to a$  (hence a constant map $u_a = p_a c_a$\inserted{$: a\to a$})
and a homotopy  $h_a : a\ \mathrm{x}\ I\to a$ from $\mathrm{id}_a$  to $u_a$ ,  i.e.
a map $h_a$ such that
\[ h_a\ (\mathrm{id}_a\,\mathrm{x}\hand{$\delta$}_o) = \mathrm{id}_a\ ,\ h_a\ (\mathrm{id}\xout{\illegible{1}}_a\,\mathrm{x}\hand{$\delta$}_1) = u_a\ . \]

We have to be cautious though,  it oddurs to me now,  not to make a
``vicious circle'',  as the \inserted{homotopy} argument used yesterday \inserted{for proving asphericity of F} makes use of the fact
that \xout{\illegible{1}}  $F\mathrm{x}I\to F$ \xout{xxx} is a weak equivalence\over{,}{.} \deleted{in order to deduce asphe-%
ricity of F.} To check that $F\mathrm{x}I\to F$ is a weak equivalence,  and even asphe-%
ric,  we have seen though that it amounts to the same to prove $a\mathrm{x}I\to a$ is a
weak equivalence for any a in A,  i.e. axI is aspheric for any a in A\unsure{?|.}
\over{x}{(}This,  then,  was seen to be a consequance of the assumption that all ele-%
ments axb are aspheric (for a,b in A)  -  but this \inserted[typed]{latter fact} is now what we want to
prove!  Thus it can't be helped,  we have to complement the condition
(T H ) above  (which I will call  (T H 1) ,  by the extra condition
\[ (2)\qquad \text{For any a in A,  axI is aspheric}\quad , \]
which,  in case I itself is in A,  appears as just a particular case of
(T 2) ,  which,\inserted{now} it seems,  we have to check directly some way or other.

19.3.

35.  I really feel I should not wait any longer with the digression on
modelizing topoi that keeps creeping in\inserted{to} my mind,  and which I keep trying
to dismiss as not to the point or not urgent or what not !  At least this
will fix some terminology,  and put the notion of an elemen\xout{x}tary modelizer
into perspective,  in terms of a \deleted{particular}\inserted{wider} class of topoi.

A topos\inserted{X} is called \uline{aspheric} if the canonical ``map'' from X to the ``final
topos'' (corresponding to a one-point space,  and \xout{xxxxx}\inserted[typed]{whose} category of sheaves
is the category (Sets)) is aspheric.  It is equivalent to say that X is
O-connected (namely non empty\inserted{,} and not decomposable non trivial\inserted[typed]{ly} into a
direct sum of two topo\"{\i}),  and that \over{foo}{\unsure{for}} any constant sheaf of groups G on
